\documentclass[11pt]{article}
\usepackage[margin=25mm]{geometry}
\usepackage{fontspec}
\setmainfont{Latin Modern Roman}
\usepackage{amsmath,amssymb,mathtools}
\usepackage{unicode-math}
\setmathfont{Latin Modern Math}
\usepackage{xurl,hyperref}
\hypersetup{colorlinks=true,urlcolor=blue}
\setcounter{secnumdepth}{0}
\setlength{\parindent}{0pt}
\setlength{\parskip}{5pt}
\setlength{\emergencystretch}{3em}
\newcommand{\sref}[2]{\hyperlink{source:#1:#2}{[S#2]}}
\newcommand{\cataloguescope}[1]{\paragraph{Recorded target.}#1}
\newcommand{\R}{\mathbb R}
\newcommand{\C}{\mathbb C}
\newcommand{\N}{\mathbb N}
\newcommand{\Z}{\mathbb Z}
\newcommand{\Q}{\mathbb Q}
\newcommand{\F}{\mathbb F}
\newcommand{\eps}{\varepsilon}
\begin{document}
\section{Equivalence of Euclidean AdS/CFT Generating Functionals}\label{statement}

\subsection{Definitions and mathematical statement}

\paragraph{Shared notation.}$\mathbb N=\{1,2,\ldots\}$; the standard number fields have their usual meanings. All additional conventions are stated locally.


\paragraph{Definitions and conventions.}
Euclidean anti-de Sitter space is hyperbolic space with a conformal boundary. A generating functional assigns a partition-function ratio to boundary data or a boundary source $h$. An exponential interaction modifies the bulk field measure, and an infrared cutoff restricts its spatial extent. Equivalence up to $h\mapsto ch$ compares the two specified functionals for all admissible sources, allowing one fixed positive normalization $c$. The thermodynamic limit removes the infrared cutoff. In the report\textquotesingle s model, the free action is $S_0(\phi)=\frac12\int_{H^d}(|\nabla\phi|^2+m^2\phi^2)\,dg$ and the cutoff interaction is $V_\Lambda(\phi)=\int_\Lambda:\!e^{\alpha\phi}\!:\,dg$. The two normalized functionals impose boundary data by a delta constraint, or pair the boundary field with a source inside an exponential, using their respective Gaussian boundary conditions.

\paragraph{Target retained from the canonical source catalogue.}
\paragraph{Catalogue wording retained for traceability.}
For the Euclidean AdS/CFT model on hyperbolic space with exponential interaction and an infrared cutoff, is the boundary-value generating functional equivalent, up to a rescaling $h\mapsto ch$ for some $c>0$, to the ordinary quantum-field-theoretic generating functional with a boundary source term? What happens to this equivalence in the thermodynamic limit?

\paragraph{Source scope and qualification.}
The primary report (PDF pages 15--16) explicitly says the finite-cutoff equivalence was established for $d=2$. It also reports triviality along some cutoff sequences in the massive exponential-interaction model; possible remedies such as holographic renormalization are left for future work. Thus the inherited sentence includes a known finite-cutoff comparison and an unresolved limiting/renormalization programme, rather than two newly verified open assertions.
\subsection{Short English statement}
Are the boundary-value and source-based Euclidean AdS/CFT generating functionals equivalent after rescaling, and does that survive removal of the cutoff?
\subsection{Research attempt}
\paragraph{Model, date and investigated scope.}
OpenAI GPT-6 Astra (configured model \texttt{gpt-6-astra}), ultra reasoning; generated 11 October 2026. Proves the finite-dimensional Gaussian source-shift identity and an explicit normalized trivial-limit countertest; preserves the source-known finite-cutoff d=2 equivalence and unresolved renormalized limit.

\paragraph{Sources and literature check.}
The complete supplied definition, including source qualifications and the dated literature assessment, was reviewed. The calculation below is an elementary restricted attempt motivated by [S2]; it does not claim novelty or a fresh literature-status determination.

\paragraph{Context and investigated scope.} The complete supplied definition, including source qualifications and the dated literature assessment, was reviewed. The calculation below is an elementary restricted attempt motivated by [S2]; it does not claim novelty or a fresh literature-status determination.

\paragraph{A finite-cutoff Gaussian identity.}
The full supplied definition says that finite-cutoff equivalence in $d=2$ was already proved and that some cutoff removals give triviality. We do not reclassify that known part as a new open result. To isolate an algebraic ingredient, let $X\in\mathbb R^N$ be centered Gaussian with positive-definite covariance $C$, let $J\in\mathbb R^N$, and let $V_\Lambda$ be a nonnegative finite-cutoff interaction. Completing the square in its Gaussian density gives
\[
 -\tfrac12x^TC^{-1}x+J^Tx
 =-\tfrac12(x-CJ)^TC^{-1}(x-CJ)+\tfrac12J^TCJ.
\]
Translation of the finite-dimensional integral therefore proves
\[
 \mathbb E[e^{J^TX}e^{-V_\Lambda(X)}]
 =e^{J^TCJ/2}\mathbb E[e^{-V_\Lambda(X+CJ)}].
\]
Both sides are finite because $0<e^{-V_\Lambda}\leq1$ and Gaussian exponential moments are finite. Dividing by the positive zero-source integral yields
\[
 \frac{Z_\Lambda(J)}{Z_\Lambda(0)}
 =e^{J^TCJ/2}
 \frac{\mathbb E[e^{-V_\Lambda(X+CJ)}]}
      {\mathbb E[e^{-V_\Lambda(X)}]}.
\]
Thus an exponential source becomes a deterministic shift, together with a Gaussian prefactor. This explains why a comparison of boundary prescriptions must track both the source normalization and the prefactor; equality of the interaction shifts alone is insufficient.

\paragraph{Why pointwise control before normalization fails.}
Even if $e^{-V_\Lambda}\leq1$, the denominator may tend to zero. For a scalar demonstration, take $X$ standard Gaussian and $V_a(x)=a x^2/2$, $a\to\infty$. Direct Gaussian integration gives
\[
 \mathbb E[e^{-V_a(X)}]=(1+a)^{-1/2}\longrightarrow0,
 \qquad
 \frac{\mathbb E[e^{JX-V_a(X)}]}{\mathbb E[e^{-V_a(X)}]}
 =\exp\left(\frac{J^2}{2(1+a)}\right)\longrightarrow1.
\]
The normalized limit has no source dependence although every finite-$a$ functional does. This quadratic example is not a replacement for the source's Wick exponential interaction; it is a countercheck on inferring nontriviality from a finite-cutoff identity.

\paragraph{Gap for the AdS/CFT problem.}
A comparison of the actual hyperbolic covariances, the boundary-to-bulk map and a fixed positive source multiplier is still needed to identify the two prescriptions. The Gaussian calculation supplies none of the infinite-dimensional ultraviolet or volume estimates. If two normalized functionals are exactly equal after a fixed source rescaling at every cutoff and both pointwise limits exist, equality passes to those limits; nontriviality does not follow. Constructing the appropriate renormalized limiting functionals and controlling the source dependence is the first missing analytical step. The reported finite-$d=2$ result and trivial cutoff sequences remain part of the source scope.

\paragraph{Verification and adversarial audit.}
Checked Gaussian completion of squares, finiteness and positivity of normalization, the explicit collapsing covariance example, and the distinction between preservation of equality and preservation of nontriviality. Model self-audit only; no independent human proof checking. Prompts A and B were applied to this canonical record and its qualified source scope. Checked Gaussian completion of squares, finiteness and positivity of normalization, the explicit collapsing covariance example, and the distinction between preservation of equality and preservation of nontriviality. Retained unresolved status and zero original-target progress; Control the actual hyperbolic boundary-to-bulk covariance comparison and normalized cutoff removal strongly enough to obtain a nontrivial generating functional.

Coordinator checked the Gaussian square completion and normalized scalar trivial-limit example, with finite-cutoff identities separated from nontriviality of a limiting field theory. This additional check was model-assisted, not independent human verification.

\paragraph{Tools, process and contributions.}
Python 3 was used for file preparation and structural checks; no numerical experiment or theorem prover was used. Contribution prompt A was followed by prompt B. GPT-6 Astra with ultra reasoning drafted the research and performed its stated self-audit; any additional coordinator contribution is identified in the verification and audit. A GPT-6-based Codex coordinator assembled the manuscript and checked repository integration; its exact variant and reasoning configuration were not exposed. No human mathematical contribution, independent human proof check or formal proof certification is supplied. Complete derivations and the consulted source references are included here; private reasoning transcripts are not published.

\paragraph{Final status and first remaining gap.}
\textbf{UNRESOLVED}. Control the actual hyperbolic boundary-to-bulk covariance comparison and normalized cutoff removal strongly enough to obtain a nontrivial generating functional.

\section{Completion Estimate}
\noindent\textbf{COMPLETION: 0\%}
The calculation is an auxiliary or restricted attempt and establishes no newly verified portion of the full source target. Control the actual hyperbolic boundary-to-bulk covariance comparison and normalized cutoff removal strongly enough to obtain a nontrivial generating functional.

\subsection{Sources}
\begin{itemize}
\item[\hypertarget{source:statement:1}{S1}] ULAM.AI and the UnsolvedMath Contributors, UnsolvedMath, catalogue entry. Catalogue curation under CC BY 4.0; definitions and English exposition added locally.
\url{https://huggingface.co/datasets/ulamai/UnsolvedMath}.
\item[\hypertarget{source:statement:2}{S2}] Recent Mathematical Developments in Quantum Field Theory (2016), Oberwolfach Reports, DOI 10.4171/owr/2016/36, pp. 2069--2122. Primary report accessed 2026-10-11; relevant PDF pages 15, 16.
\url{https://doi.org/10.4171/owr/2016/36}.
\item[\hypertarget{source:statement:3}{S3}] Gottschalk--Thaler, Euclidean AdS/CFT (catalogue-listed context; not independently consulted in this preparation).
\url{https://arxiv.org/abs/math-ph/0611006}.
\item[\hypertarget{source:statement:4}{S4}] Gottschalk--Thaler, triviality (catalogue-listed context; not independently consulted in this preparation).
\url{https://doi.org/10.1007/s00220-013-1796-9}.
\end{itemize}
\label{end:statement}
\end{document}
